
\section{Original occurrences and common sample assignments}
\label{sec:peps_original_gate_sample_reindexing}

Following \cite[proof of Theorem~5.2]{OpenAI2026PolynomialPEPS}, let
$\mathcal W:a\to b$ be a chronological source circuit on an arbitrary
party set $P$. Write $G$ for its original gate occurrences with nonempty
participant sets $C_g$, and $\Xi_g$ for their finite branch sets. Repeated
uses of a gate are distinct occurrences. Write $\mathcal E(C_g)$ for the
unordered pairs of distinct participants and
$\mathcal S=\bigsqcup_{g\in G}\mathcal E(C_g)$ for the shared sample positions.
Let $\mathcal S_p$ consist of positions whose pair contains $p$.
The original ordered full source list has length $N$, and the original
ordered incident source list at $p$ has length $n_p$.
For $r\in\N$, write $[r]=\{0,\ldots,r-1\}$.
The canonical bijections of source positions are
$\tau:[N]\to\mathcal S$ and $\kappa_p:[n_p]\to\mathcal S_p$.
% Source: 04-compression.tex, lines 342--417 and 565--588.

\begin{definition}[A participant root at each original gate]
    \label{def:peps_source_gate_root}
    \lean{TNLean.PEPS.PairEffect.SourceCircuit.sourceGateRoot}
    \leanok
    \uses{def:peps_touched_labels}
    In \cite[proof of Theorem~5.2]{OpenAI2026PolynomialPEPS}, choose
    $r_g\in C_g$ for every $g\in G$. The resulting assignment
    $r:G\to P$ specifies the root of the star sharing the gate label.
    Each $C_g$ is nonempty by the definition of $G$. No nonemptiness
    assumption on $P$ is imposed; when $G$ is empty the assignment
    is the unique empty function.
    % Source: 04-compression.tex, lines 565--579.
\end{definition}

\begin{theorem}[The root is an original participant]
    \label{thm:peps_source_gate_root_mem}
    \lean{TNLean.PEPS.PairEffect.SourceCircuit.sourceGateRoot_mem}
    \leanok
    \uses{def:peps_source_gate_root}
    For the star roots in \cite[proof of Theorem~5.2]{OpenAI2026PolynomialPEPS},
    the selected root at every original occurrence $g\in G$ satisfies
    $r_g\in C_g$, and the party set is arbitrary.
    % Source: 04-compression.tex, lines 565--579.
\end{theorem}

\begin{proof}
    \leanok
    \uses{def:peps_source_gate_root}
    By definition of $G$, each participant set $C_g$ is nonempty and $r_g$ is chosen from $C_g$.
\end{proof}

\begin{theorem}[Incident labels after reindexing a full choice]
    \label{thm:peps_reindexed_incident_gate_labels}
    \lean{TNLean.PEPS.PairEffect.SourceCircuit.originalGateLabelsAtOfAllChoices_eq_restrictGates}
    \leanok
    \uses{thm:peps_incident_gate_labels,def:peps_touched_labels}
    For the gate-label identification in
    \cite[proof of Theorem~5.2]{OpenAI2026PolynomialPEPS}, let $e$ be the
    canonical bijection from all-party chronological
    choices to $\prod_{g\in G}\Xi_g$. For any original label function
    $\eta\in\prod_{g\in G}\Xi_g$ and any $p\in P$, extracting the
    incident labels of $e^{-1}\eta$ gives precisely
    $\eta|_{\{g:p\in C_g\}}$. This holds for empty incident sets and,
    when $\eta$ is given, empty branch sets.
    % Source: 04-compression.tex, lines 351--381 and 565--579.
\end{theorem}

\begin{proof}
    \leanok
    \uses{thm:peps_incident_gate_labels,def:peps_touched_labels}
    At an incident occurrence $g$, the extracted value is
    $(e(e^{-1}\eta))_g=\eta_g$. Equality at every incident occurrence
    gives the assertion.
\end{proof}

\begin{theorem}[Incident samples after reindexing a full assignment]
    \label{thm:peps_reindexed_incident_samples}
    \lean{TNLean.PEPS.PairEffect.SourceCircuit.restrict_reindexedSourceSampleAssignment_eq}
    \leanok
    \uses{def:peps_source_sample_assignment_indices,def:peps_local_source_sample_indices}
    For the local sample values in
    \cite[proof of Theorem~5.2]{OpenAI2026PolynomialPEPS}, let $k\in\N$ and
    let $j:\mathcal S\to[k]$. The corresponding
    assignment on the full ordered list is $\widetilde j=j\circ\tau$.
    If $\iota_p:[n_p]\to[N]$ reads the original positions of the
    incident slots, then
    $\widetilde j\circ\iota_p=(j|_{\mathcal S_p})\circ\kappa_p$.
    No finiteness assumption on $P$ is required, and this holds for $k=0$
    and empty incident lists for any given $j$.
    % Source: 04-compression.tex, lines 342--417 and 568--579.
\end{theorem}

\begin{proof}
    \leanok
    \uses{def:peps_source_sample_assignment_indices,def:peps_local_source_sample_indices}
    The inverse assignment reads $j$ at the sample position of each
    original slot. Since $\tau(\iota_p(i))=\kappa_p(i)$, both sides
    at $i$ equal $j(\kappa_p(i))$.
\end{proof}

\begin{theorem}[Reindexing choices and common samples together]
    \label{thm:peps_joint_gate_sample_sum}
    \lean{TNLean.PEPS.PairEffect.SourceCircuit.sum_choiceSamples_eq_sum_originalGateSamples}
    \leanok
    \uses{def:peps_touched_labels,def:peps_source_sample_assignment_indices}
    In \cite[proof of Theorem~5.2]{OpenAI2026PolynomialPEPS}, let $P$ be
    arbitrary, $k\in\N$, and let $I,J$ be arbitrary index sets.
    Let $\Xi_{\rm all}$ be the all-party chronological choice set,
    $c_\xi$ its coefficient, and $a_g:\Xi_g\to\C$ the original gate
    coefficient. For any matrix family
    $K:\Xi_{\rm all}\times\Xi_{\rm all}\times[k]^{[N]}\to\C^{I\times J}$,
    one has
    \begin{align}
        \sum_{\xi,\zeta}c_\xi\overline{c_\zeta}
        \sum_{\ell:[N]\to[k]}K(\xi,\zeta,\ell)
         & =\sum_{\eta,\theta\in\prod_{g\in G}\Xi_g}
        \left(\prod_{g\in G}a_g(\eta_g)\overline{a_g(\theta_g)}\right)
        \sum_{j:\mathcal S\to[k]}K(e^{-1}\eta,e^{-1}\theta,j\circ\tau).
        \label{eq:peps_joint_gate_sample_sum}
    \end{align}
    The two branch functions are independent; the sample assignment
    is shared. No finiteness of $I,J$ or $P$ is assumed, and the
    empty-participant scalar is outside the coefficient product.
    % Source: 04-compression.tex, lines 351--381 and 565--579.
\end{theorem}

\begin{proof}
    \leanok
    \uses{thm:peps_mixed_gate_label_sum,thm:peps_source_sample_assignment_sum}
    Reindexing the mixed chronological choice sum independently in its
    two arguments and reindexing the common sample assignment once by
    $\tau$ in every resulting summand gives
    (\ref{eq:peps_joint_gate_sample_sum}).
\end{proof}

\section{Ordered gate labels in a distributed density contraction}
\label{sec:peps_ordered_gate_density_algebra}

\begin{theorem}[Labels on the full selected gate set]
    \label{thm:peps_full_selected_gate_label_sum}
    \lean{TNLean.PEPS.Approximation.DistributedOperatorContraction.sum_eq_sum_gateLabels_univ}
    \leanok
    \uses{def:compression_operator_gate_labels}
    For the occurrence identification in
    \cite[proof of Theorem~5.2]{OpenAI2026PolynomialPEPS}, let $G$ be finite
    and let $\Lambda_g$ be finite label sets. Let
    $\widetilde G$ denote the full selected gate set, with its canonical
    bijection $u:\widetilde G\to G$. For arbitrary index sets $I,J$
    and $K:(\prod_{g\in G}\Lambda_g)\to\C^{I\times J}$,
    \begin{align}
        \sum_{\eta\in\prod_{g\in G}\Lambda_g}K(\eta)
         & =\sum_{\xi\in\prod_{h\in\widetilde G}\Lambda_{u(h)}}
        K(\xi\circ u^{-1}).
        \label{eq:peps_full_selected_gate_label_sum}
    \end{align}
    Empty gate and label sets are included.
    % Source: 04-compression.tex, lines 565--579.
\end{theorem}

\begin{proof}
    \leanok
    Precomposition with $u^{-1}$ is a bijection of the dependent label
    assignments, and reindexing the finite matrix sum by this bijection
    yields the stated sum identity.
\end{proof}

For the next three statements, let $G$ be finite and let $K_g,B_g$ be
finite ket and bra label sets. Write $\mathcal K=\prod_g K_g$ and
$\mathcal B=\prod_g B_g$. For $z\in\prod_g(K_g\times B_g)$, write
$z^K_g=(z_g)_1$ and $z^B_g=(z_g)_2$. The matrix indices $I,J$ are arbitrary.

\begin{theorem}[Pairing the ket and bra label sums]
    \label{thm:peps_ordered_gate_label_pair_sum}
    \lean{TNLean.PEPS.Approximation.sum_gateLabels_eq_sum_orderedGateLabelPairs}
    \leanok
    For the shared ket and bra label in
    \cite[proof of Theorem~5.2]{OpenAI2026PolynomialPEPS}, any matrix family
    $F:\mathcal K\times\mathcal B\to\C^{I\times J}$ satisfies
    \begin{align}
        \sum_{\eta\in\mathcal K}\sum_{\theta\in\mathcal B}F(\eta,\theta)
         & =\sum_{z\in\prod_g(K_g\times B_g)}F(z^K,z^B).
        \label{eq:peps_ordered_gate_label_pair_sum}
    \end{align}
    The pairs are ordered, and empty gate and label sets are included.
    % Source: 04-compression.tex, lines 565--579.
\end{theorem}

\begin{proof}
    \leanok
    Replacing the iterated sum by the sum over $\mathcal K\times\mathcal B$
    and reindexing along the bijection
    $(\eta,\theta)\mapsto((\eta_g,\theta_g))_g$ gives
    (\ref{eq:peps_ordered_gate_label_pair_sum}).
\end{proof}

\begin{theorem}[Pairing the weighted ket and bra sums]
    \label{thm:peps_ordered_gate_label_pair_weighted_sum}
    \lean{TNLean.PEPS.Approximation.sum_weighted_gateLabels_eq_sum_orderedGateLabelPairs}
    \leanok
    For the ordered coefficient sum in
    \cite[proof of Theorem~5.2]{OpenAI2026PolynomialPEPS}, arbitrary
    $a_g:K_g\to\C$, $b_g:B_g\to\C$ and the matrix family $F$ satisfy
    \begin{align}
        \sum_{\eta,\theta}
        \left(\prod_g a_g(\eta_g)\overline{b_g(\theta_g)}\right)F(\eta,\theta)
         & =\sum_z\left(\prod_g a_g(z^K_g)\overline{b_g(z^B_g)}\right)
        F(z^K,z^B).
        \label{eq:peps_ordered_gate_label_pair_weighted_sum}
    \end{align}
    No nonvanishing or normalization is assumed.
    % Source: 04-compression.tex, lines 565--579.
\end{theorem}

\begin{proof}
    \leanok
    \uses{thm:peps_ordered_gate_label_pair_sum}
    Applying (\ref{eq:peps_ordered_gate_label_pair_sum}) to
    $(\eta,\theta)\mapsto(\prod_g a_g(\eta_g)\overline{b_g(\theta_g)})F(\eta,\theta)$
    gives the weighted-sum identity.
\end{proof}

\begin{theorem}[Placing gate weights inside the ordered coefficient product]
    \label{thm:peps_ordered_gate_label_averaging}
    \lean{TNLean.PEPS.Approximation.weighted_sum_gateLabels_eq_exterior_smul_orderedGateLabelPairs}
    \leanok
    To place the averaging factors in the gate coefficients as in
    \cite[proof of Theorem~5.2]{OpenAI2026PolynomialPEPS}, let $w_g\in\C$ be
    arbitrary gate weights and $\gamma\in\C$ an exterior scalar. Then
    \begin{align}
        \left(\prod_g w_g\right)\gamma
        \sum_{\eta,\theta}\left(\prod_g a_g(\eta_g)\overline{b_g(\theta_g)}\right)
        F(\eta,\theta)
         & =\gamma\sum_z
        \left(\prod_g w_g a_g(z^K_g)\overline{b_g(z^B_g)}\right)F(z^K,z^B).
        \label{eq:peps_ordered_gate_label_averaging}
    \end{align}
    All weights may vanish, and all empty finite-sum cases are included.
    % Source: 04-compression.tex, lines 568--588.
\end{theorem}

\begin{proof}
    \leanok
    \uses{thm:peps_ordered_gate_label_pair_weighted_sum}
    By (\ref{eq:peps_ordered_gate_label_pair_weighted_sum}), distributing
    $\prod_g w_g$ over the finite sum and multiplying the two products
    term by term yields the averaging identity, with
    the factor $\gamma$ remaining outside throughout.
\end{proof}

\begin{theorem}[The distributed density contraction in original gate labels]
    \label{thm:peps_density_pair_univ_double_sum}
    \lean{TNLean.PEPS.Approximation.DistributedOperatorContraction.contractedOperator_densityPair_univ_eq_double_sum}
    \leanok
    \uses{def:compression_operator_operator,def:compression_operator_density_coefficient}
    In \cite[proof of Theorem~5.2]{OpenAI2026PolynomialPEPS}, let $P,G$ be
    finite, and let $C_g\subseteq P$ be arbitrary participant
    sets with chosen roots $r_g\in C_g$. Let $K_g,B_g$ be finite label
    sets, $k\in\N$, $a_g:K_g\to\C$, $b_g:B_g\to\C$,
    $w_g\in\C$ and $\gamma\in\C$.
    Let $O_p,I_p$ be arbitrary matrix index sets, without finiteness
    assumptions. Put $\mathcal S=\bigsqcup_g\mathcal E(C_g)$ and let
    $\mathcal S_p$ consist of its positions incident to $p$.
    Let
    $T_p:\prod_{g:p\in C_g}K_g\times\prod_{g:p\in C_g}B_g
        \times[k]^{\mathcal S_p}\to\C^{O_p\times I_p}$
    be a family of matrices.
    Write $\mathcal C(C,r,K\times B,k,c,T)$ for the distributed
    contraction with ordered gate-pair labels, roots $r_g$, common
    sample positions, coefficient
    $c_g(\alpha,\beta)=w_g a_g(\alpha)\overline{b_g(\beta)}$, and local
    matrix $T_p$ receiving the two coordinates independently. Then
    \begin{align}
        \gamma\mathcal C(C,r,K\times B,k,c,T)
         & =\left(\prod_g w_g\right)\gamma
        \sum_{\eta\in\mathcal K}\sum_{\theta\in\mathcal B}
        \left(\prod_g a_g(\eta_g)\overline{b_g(\theta_g)}\right)
        \sum_{j:\mathcal S\to[k]}
        \bigotimes_{p\in P}T_p(\eta|_p,\theta|_p,j|_p).
        \label{eq:peps_density_pair_univ_double_sum}
    \end{align}
    The tensor product has indices $\prod_pO_p$ and $\prod_pI_p$.
    Empty families, empty labels, $k=0$ and zero weights are included
    whenever the stated root data exist.
    % Source: 04-compression.tex, lines 565--588.
\end{theorem}

\begin{proof}
    \leanok
    \uses{thm:compression_operator_density_pair,
        thm:peps_full_selected_gate_label_sum,thm:peps_ordered_gate_label_averaging}
    Evaluating the density-pair contraction as a sum over labels on the
    full selected gate set and reindexing those labels by
    (\ref{eq:peps_full_selected_gate_label_sum}) leaves the coefficient
    product unchanged under the same occurrence bijection, while each local
    matrix receives the same incident coordinates and restricted samples.
    Applying (\ref{eq:peps_ordered_gate_label_averaging}) in reverse then
    yields (\ref{eq:peps_density_pair_univ_double_sum}).
\end{proof}

\section{The sampled physical density as a distributed contraction}
\label{sec:peps_actual_sampled_distributed_contraction}

Let $P$ be finite, let $d:P\to\N$, and let an ordered list record every
party exactly once. Let $h_d$ be its original physical register list,
with space $\C^{d_p}$ at $p$. Let $c$ be an arbitrary private output
list and $\mathcal W:a\to h_dc$ a chronological circuit with recorded
product input $\psi$.
At every original source $e$, retain the original endpoint dimensions
$u_e,v_e$ and $r_e=\min(u_e,v_e)$. Fix arbitrary real coefficient
families and complex endpoint matrices of sizes $u_e\times r_e$ and
$v_e\times r_e$, with their original branch indices, sample number
$k\in\N$ and original ordered-pair arrays $\omega$.
At each all-party grouped owner, choose finite-index orthonormal bases
of the selected source memory and private output memory, together with
finite-index orthonormal bases of the original input
memory and original private output memory. Use the original labelled
physical basis and the physical/private output bases constructed from
these data, as in Definition~\ref{def:peps_incident_sampled_physical_factor}.
Write $\widehat\rho_{\mathrm{phys},\psi}(\widehat X)$ for the original
physical operator with the sampled source matrices, and
$\gamma=\lambda_W\overline{\lambda_W}$ for the squared modulus of its
original empty-participant gate scalar.
Let $a_g$ denote the original branch coefficient at $g\in G$.
For incident gate labels $\alpha,\beta$ and a sample assignment
$s:\mathcal S_p\to[k]$, put
\begin{align}
    T_p(\alpha,\beta,s)
     & =M^{\rm inc}_{p,\alpha,\beta,s\circ\kappa_p,\omega}.
    \label{eq:peps_distributed_incident_matrix}
\end{align}
Thus every matrix retains the original physical indices $[d_p]$, and
its sample argument is read at the original ordered incident slots.

\begin{theorem}[The sampled physical sum in original gate and sample positions]
    \label{thm:peps_original_gate_sample_physical_sum}
    \lean{TNLean.PEPS.PairEffect.SourceCircuit.physicalSourceReplacedDensity_pure_sampled_eq_scalar_smul_originalGateSampleSum}
    \leanok
    \uses{def:peps_incident_sampled_physical_factor,def:peps_source_sample_assignment_indices}
    In \cite[proof of Theorem~5.2]{OpenAI2026PolynomialPEPS}, the data just
    specified satisfy
    \begin{align}
        \widehat\rho_{\mathrm{phys},\psi}(\widehat X)
         & =(k^{-1})^N\gamma
        \sum_{\eta,\theta\in\prod_{g\in G}\Xi_g}
        \left(\prod_{g\in G}a_g(\eta_g)\overline{a_g(\theta_g)}\right)
        \sum_{j:\mathcal S\to[k]}
        \bigotimes_{p\in P}T_p(\eta|_p,\theta|_p,j|_p).
        \label{eq:peps_original_gate_sample_physical_sum}
    \end{align}
    The two label functions are independent and the sample assignment
    is shared. The matrices use the same original dimensions, bases
    and arrays $\omega$. Empty party and source sets, zero dimensions,
    vanishing coefficients and $k=0$ are included with $0^{-1}=0$ and
    zeroth powers equal to $1$.
    % Source: 04-compression.tex, lines 351--381 and 565--588.
\end{theorem}

\begin{proof}
    \leanok
    \uses{thm:peps_incident_pure_sampled_physical_sum,
        thm:peps_joint_gate_sample_sum,thm:peps_reindexed_incident_gate_labels,
        thm:peps_reindexed_incident_samples}
    Starting from the original sum in incident factors,
    (\ref{eq:peps_joint_gate_sample_sum}) reindexes its two branch choices
    and common sample assignment. The incident-label and incident-sample
    identities turn each matrix into
    (\ref{eq:peps_distributed_incident_matrix}) evaluated at
    $\eta|_p,\theta|_p,j|_p$, while the original scalar and arrays remain
    unchanged, giving (\ref{eq:peps_original_gate_sample_physical_sum}).
\end{proof}

\begin{theorem}[The exact distributed contraction of the sampled physical density]
    \label{thm:peps_actual_sampled_distributed_contraction}
    \lean{TNLean.PEPS.PairEffect.SourceCircuit.physicalSourceReplacedDensity_pure_sampled_eq_exterior_smul_contractedOperator}
    \leanok
    \uses{def:compression_operator_operator,def:compression_operator_density_coefficient,
        def:peps_incident_sampled_physical_factor,def:peps_source_gate_root}
    In \cite[proof of Theorem~5.2]{OpenAI2026PolynomialPEPS}, retain the
    original physical data, sample arrays, matrices $T_p$ and scalar
    $\gamma$ specified in this section. Choose the
    recorded participant roots $r_g$ and define
    \begin{align}
        w_g               & =(k^{-1})^{|\mathcal E(C_g)|},         \\
        c_g(\alpha,\beta) & =w_g a_g(\alpha)\overline{a_g(\beta)}.
        \label{eq:peps_actual_gate_weight}
    \end{align}
    The exact physical operator satisfies
    \begin{align}
        \widehat\rho_{\mathrm{phys},\psi}(\widehat X)
         & =\gamma\mathcal C(C,r,\Xi\times\Xi,k,c,T).
        \label{eq:peps_actual_sampled_distributed_contraction}
    \end{align}
    The star at each original gate shares the entire ordered ket and
    bra pair. Each original source occurrence shares one sample index
    between its two original participants. Its averaging factor occurs
    once in that gate's coefficient at the root. The scalar $\gamma$
    remains outside the contraction, and all empty and zero cases of
    Theorem~\ref{thm:peps_original_gate_sample_physical_sum} are included.
    % Source: 04-compression.tex, lines 565--588.
\end{theorem}

\begin{proof}
    \leanok
    \uses{thm:peps_original_gate_sample_physical_sum,
        thm:peps_source_gate_root_mem,thm:peps_density_pair_univ_double_sum,
        thm:peps_original_source_sample_averaging}
    Evaluating the distributed contraction by
    (\ref{eq:peps_density_pair_univ_double_sum}) with the original
    gate coefficients for both the ket and bra, where the selected roots
    are original participants, and replacing $\prod_g w_g$ by $(k^{-1})^N$
    via (\ref{eq:peps_original_source_sample_averaging}) yields
    (\ref{eq:peps_original_gate_sample_physical_sum}) and hence
    (\ref{eq:peps_actual_sampled_distributed_contraction}).
\end{proof}
